% CP-VI-0042
% target_chapter: VI/30 — Dimension of Schemes
% source_unit_id: IMP-DL-A5CEE6E541-P002
% source: Downloads/theory-of-algebraic-geometry-8.tex L340-L747
% solution_provenance: SOURCE_IMPLICIT_SOLUTION

\subsection*{Problem 2: Dimension and Codimension for Varieties}

\subsection*{Statement}

Let \(X\) be a variety over \(\operatorname{Spec} k\). Show:

\begin{enumerate}[label=(\alph*)]
	\item For any closed point \(p\in X\),
	\[
	\dim X=\dim \mathcal O_{X,p}.
	\]
	
	\item If \(Y\subseteq X\) is a closed subset of \(X\), define
	\[
	\operatorname{codim}(Y,X)
	=
	\inf_{Z\subseteq Y}\operatorname{codim}(Z,X),
	\]
	where the infimum is taken over all closed irreducible subsets \(Z\) of \(Y\). Show that
	\[
	\operatorname{codim}(Y,X)
	=
	\inf\{\dim \mathcal O_{X,y}\mid y\in Y\}.
	\]
	
	\item If \(Y\subseteq X\) is closed, then
	\[
	\dim Y+\operatorname{codim}(Y,X)=\dim X.
	\]
	
	\item If \(U\subseteq X\) is a nonempty open subset, then
	\[
	\dim U=\dim X.
	\]
	
	\item \(\dim X\) coincides with the transcendence degree of \(K(X)\) over \(k\):
	\[
	\dim X=\trdeg_k K(X).
	\]
\end{enumerate}

We may use the following facts from commutative algebra.

If \(B\) is a finitely generated \(k\)-algebra and a domain, and if \(\mathfrak p\subseteq B\) is a prime ideal, then
\[
\operatorname{height}(\mathfrak p)+\dim B/\mathfrak p=\dim B.
\]

If \(K\) is the field of fractions of \(B\), then
\[
\dim B=\trdeg_k K.
\]

\subsection*{General Setup}

Since \(X\) is a variety over \(k\), we assume that \(X\) is an integral scheme of finite type over \(k\).

Let
\[
K(X)
\]
denote the function field of \(X\).

If
\[
U=\operatorname{Spec} B\subseteq X
\]
is a nonempty affine open subset, then \(B\) is a finitely generated \(k\)-algebra and a domain. Moreover,
\[
\operatorname{Frac}(B)=K(X).
\]

By the given commutative algebra fact,
\[
\dim B=\trdeg_k \operatorname{Frac}(B).
\]
Since \(\operatorname{Frac}(B)=K(X)\), this gives
\[
\dim B=\trdeg_k K(X).
\]

This observation will be used several times.

\subsection{Part (a): Closed Points and Local Dimension}

Let \(p\in X\) be a closed point. We need to prove that
\[
\dim X=\dim \mathcal O_{X,p}.
\]

Choose an affine open neighbourhood
\[
U=\operatorname{Spec} B\subseteq X
\]
of \(p\). The point \(p\) corresponds to a prime ideal \(\mathfrak m\subseteq B\). Since \(p\) is closed in \(X\), and hence closed in the affine open neighbourhood \(U\), the ideal \(\mathfrak m\) is maximal.

Thus
\[
p\leftrightarrow \mathfrak m\in \operatorname{Spec} B.
\]

Because \(B\) is a finitely generated \(k\)-algebra and \(\mathfrak m\) is maximal, the quotient \(B/\mathfrak m\) is a finite field extension of \(k\). In particular,
\[
\dim B/\mathfrak m=0.
\]

Using the given dimension formula for finitely generated \(k\)-algebras, we have
\[
\operatorname{height}(\mathfrak m)+\dim B/\mathfrak m=\dim B.
\]
Since \(\dim B/\mathfrak m=0\), it follows that
\[
\operatorname{height}(\mathfrak m)=\dim B.
\]

On the other hand,
\[
\mathcal O_{X,p}\cong B_{\mathfrak m}.
\]
The dimension of the local ring \(B_{\mathfrak m}\) is the height of \(\mathfrak m\):
\[
\dim B_{\mathfrak m}=\operatorname{height}(\mathfrak m).
\]
Hence
\[
\dim \mathcal O_{X,p}
=
\operatorname{height}(\mathfrak m).
\]

Therefore,
\[
\dim \mathcal O_{X,p}=\dim B.
\]

Since \(U=\operatorname{Spec} B\) is a nonempty affine open subset of the variety \(X\), we have
\[
\dim B=\dim X.
\]
Thus
\[
\boxed{\dim \mathcal O_{X,p}=\dim X.}
\]

\subsection{Part (b): Codimension of a Closed Subset}

Let \(Y\subseteq X\) be a closed subset. We are given the definition
\[
\operatorname{codim}(Y,X)
=
\inf_{Z\subseteq Y}\operatorname{codim}(Z,X),
\]
where \(Z\) runs over the closed irreducible subsets of \(Y\).

By Problem 1, every closed irreducible subset \(Z\subseteq Y\) has a unique generic point \(\eta_Z\in Z\) such that
\[
\overline{\{\eta_Z\}}=Z.
\]

Conversely, if \(y\in Y\), then since \(Y\) is closed in \(X\), we have
\[
\overline{\{y\}}\subseteq Y.
\]
The closure \(\overline{\{y\}}\) is always irreducible, because it is the closure of a single point.

Thus the closed irreducible subsets of \(Y\) are exactly the closures
\[
\overline{\{y\}},
\qquad y\in Y.
\]

Now fix \(y\in Y\). Choose an affine open neighbourhood
\[
U=\operatorname{Spec} B\subseteq X
\]
of \(y\). Let \(y\) correspond to the prime ideal
\[
\mathfrak p\subseteq B.
\]
Then
\[
\overline{\{y\}}\cap U=V(\mathfrak p).
\]

The codimension of the irreducible closed subset \(\overline{\{y\}}\) in \(X\) is locally the height of \(\mathfrak p\):
\[
\operatorname{codim}(\overline{\{y\}},X)
=
\operatorname{height}(\mathfrak p).
\]

Also,
\[
\mathcal O_{X,y}\cong B_{\mathfrak p}.
\]
Therefore
\[
\dim \mathcal O_{X,y}
=
\dim B_{\mathfrak p}
=
\operatorname{height}(\mathfrak p).
\]

Hence
\[
\operatorname{codim}(\overline{\{y\}},X)=\dim \mathcal O_{X,y}.
\]

Taking the infimum over all \(y\in Y\), we obtain
\[
\operatorname{codim}(Y,X)
=
\inf_{y\in Y}\operatorname{codim}(\overline{\{y\}},X)
=
\inf_{y\in Y}\dim \mathcal O_{X,y}.
\]

Therefore
\[
\boxed{
	\operatorname{codim}(Y,X)
	=
	\inf\{\dim \mathcal O_{X,y}\mid y\in Y\}.
}
\]

\subsection{Part (c): Dimension Plus Codimension}

Let \(Y\subseteq X\) be a closed subset. We need to show that
\[
\dim Y+\operatorname{codim}(Y,X)=\dim X.
\]

Let \(Z\subseteq Y\) be a closed irreducible subset. Let \(\eta\in Z\) be its generic point, so that
\[
Z=\overline{\{\eta\}}.
\]

Choose an affine open subset
\[
U=\operatorname{Spec} B\subseteq X
\]
meeting \(Z\). The point \(\eta\) corresponds to a prime ideal
\[
\mathfrak p\subseteq B.
\]
Then
\[
Z\cap U=V(\mathfrak p).
\]

The codimension of \(Z\) in \(X\) is
\[
\operatorname{codim}(Z,X)=\operatorname{height}(\mathfrak p).
\]

The dimension of \(Z\) is
\[
\dim Z=\dim B/\mathfrak p.
\]

Using the given formula,
\[
\operatorname{height}(\mathfrak p)+\dim B/\mathfrak p=\dim B.
\]
Therefore
\[
\operatorname{codim}(Z,X)+\dim Z=\dim B.
\]

Since \(U=\operatorname{Spec} B\) is a nonempty affine open subset of \(X\), we have
\[
\dim B=\dim X.
\]
Hence
\[
\operatorname{codim}(Z,X)+\dim Z=\dim X.
\]

Thus
\[
\operatorname{codim}(Z,X)=\dim X-\dim Z.
\]

Now take the infimum over all closed irreducible subsets \(Z\subseteq Y\):
\[
\operatorname{codim}(Y,X)
=
\inf_{Z\subseteq Y}\operatorname{codim}(Z,X).
\]
Using the formula above,
\[
\operatorname{codim}(Y,X)
=
\inf_{Z\subseteq Y}(\dim X-\dim Z).
\]

Since \(\dim X\) is fixed, this becomes
\[
\operatorname{codim}(Y,X)
=
\dim X-\sup_{Z\subseteq Y}\dim Z.
\]

But by definition,
\[
\dim Y
=
\sup\{\dim Z\mid Z\subseteq Y \text{ closed irreducible}\}.
\]

Therefore
\[
\operatorname{codim}(Y,X)=\dim X-\dim Y.
\]

Rearranging gives
\[
\boxed{\dim Y+\operatorname{codim}(Y,X)=\dim X.}
\]

\subsection{Part (d): Nonempty Open Subsets Have the Same Dimension}

Let \(U\subseteq X\) be a nonempty open subset. We need to show that
\[
\dim U=\dim X.
\]

Since \(X\) is integral, every nonempty open subset of \(X\) is also integral. Therefore \(U\) is integral.

Since \(X\) is of finite type over \(k\), the open subscheme \(U\) is also of finite type over \(k\). Hence \(U\) is again a variety over \(k\).

Because \(U\) is a nonempty open subset of the integral scheme \(X\), \(U\) and \(X\) have the same function field:
\[
K(U)=K(X).
\]

Using Part (e), which states that the dimension of a variety is the transcendence degree of its function field, we get
\[
\dim U=\trdeg_k K(U)
\]
and
\[
\dim X=\trdeg_k K(X).
\]

Since
\[
K(U)=K(X),
\]
we conclude that
\[
\dim U=\dim X.
\]

Hence
\[
\boxed{\dim U=\dim X.}
\]

\subsection{Part (e): Dimension Equals Transcendence Degree}

We prove that
\[
\dim X=\trdeg_k K(X).
\]

Choose a nonempty affine open subset
\[
U=\operatorname{Spec} B\subseteq X.
\]
Since \(X\) is a variety over \(k\), the ring \(B\) is a finitely generated \(k\)-algebra and a domain.

The function field of \(X\) is the fraction field of \(B\):
\[
K(X)=\operatorname{Frac}(B).
\]

By the commutative algebra fact given in the question,
\[
\dim B=\trdeg_k \operatorname{Frac}(B).
\]
Therefore
\[
\dim B=\trdeg_k K(X).
\]

Since \(U=\operatorname{Spec} B\) is a nonempty affine open subset of the irreducible finite-type scheme \(X\), it has the same dimension as \(X\):
\[
\dim U=\dim X.
\]
But
\[
\dim U=\dim B.
\]
Thus
\[
\dim X=\dim B=\trdeg_k K(X).
\]

Therefore
\[
\boxed{\dim X=\trdeg_k K(X).}
\]

\newpage
